\section{Only two possible crossing exponents}
\label{sec:crossing-exponent-gap}

The tree-price dichotomy sharpens the crossing formulation:
on the nonprice branch the crossing exponent is its elementary
lower bound, rather than an unspecified improvement below \(\sqrt2\).
This conclusion still leaves the branch undecided.
Over \(\mathbb C\), Liu's unrefereed preprint already derives the sharp
\(\sqrt2\) Hurwitz exponent from the absence of a planar strict
shortening \cite[p.~3, \emph{Two crossing cuts}]{Liu2026noncommutative}.
The exponent-one branch below uses the tree-price dichotomy.

\begin{theorem}[A gap in the crossing exponent]
\label{thm:crossing-exponent-dichotomy}
Fix a field \(\kk\). Assume the price-space theorem
\cite[Proposition C.3, p.~28]{Anthropic2026} and the cover duality
of Theorem~\ref{thm:coverduality}. Then the limit
\[
 \rho_\kk=\lim_{q\to\infty}\frac{\log P_\kk(q)}{\log q}
\]
exists and satisfies
\begin{equation}\label{eq:crossing-exponent-dichotomy}
 \rho_\kk=
 \begin{cases}
  1,&\text{if the unit square is not a price over }\kk,\\
  \sqrt2,&\text{if the unit square is a price over }\kk.
 \end{cases}
\end{equation}
In particular, any fixed exponent improvement below \(\sqrt2\)
along an unbounded sequence of dimensions implies
\(P_\kk(q)=q^{1+o(1)}\) for all integer \(q\to\infty\).
\end{theorem}
\begin{proof}
Suppose first that the square is not a price.
By Theorem~\ref{thm:tree-price-dichotomy}, every finite price
is a tree pseudometric. Thus, for every relaxed spectral point
\(\phi\) on four parties, its wire metric satisfies
\[
 d_\phi(1,3)+d_\phi(2,4)
 \leq
 \max\{d_\phi(1,2)+d_\phi(3,4),
        d_\phi(1,4)+d_\phi(2,3)\}.
\]
Put
\[
 C=\wire{2}{1}{3}\wire{2}{2}{4},\qquad
 A_0=\wire{2}{1}{2}\wire{2}{3}{4},\qquad
 A_1=\wire{2}{1}{4}\wire{2}{2}{3}.
\]
Multiplicativity converts the preceding inequality into
\[
 \phi(C)\leq\max\{\phi(A_0),\phi(A_1)\}.
\]
For every \(\varepsilon>0\), cover duality therefore supplies,
for all sufficiently large \(N\), a cover of \(C^N\) by at most
\(2^{\varepsilon N}\) restrictions of products of \(N\) factors
chosen from \(A_0,A_1\).

A resource product with \(j\) copies of \(A_0\) and \(N-j\)
copies of \(A_1\) has ranks
\[
 2^{2(N-j)}\quad\hbox{across }12\mid34,
 \qquad
 2^{2j}\quad\hbox{across }14\mid23.
\]
Its smaller cut rank is at most \(2^N\). Choose that cut for
each cover term, retain the same bound after its local restriction,
and group terms according to the chosen cut. Matrix-rank
subadditivity gives
\begin{equation}\label{eq:tree-crossing-cover-cost}
 P_\kk(2^N)\leq2^{(1+\varepsilon)N}.
\end{equation}
There is no additional \(2^N\) choice cost here: the cover
duality already bounds the number of selected product terms by
\(2^{\varepsilon N}\).

For arbitrary \(q\), restrict from \(2^N\) with
\(N=\lceil\log_2q\rceil\). Then
\[
 P_\kk(q)\leq
 2^{(1+\varepsilon)\lceil\log_2q\rceil}.
\]
Taking the upper limit and then \(\varepsilon\downarrow0\),
and using \(P_\kk(q)\geq q\) from Lemma~\ref{lem:pt-upper},
proves that the exponent is one.

If the square is a price, Theorem~\ref{thm:square-pt-exponent}
gives exponent \(\sqrt2\). These two alternatives exhaust the
possibilities. A fixed improvement below \(\sqrt2\) on an unbounded
sequence excludes the price branch, and hence forces the first
branch of \eqref{eq:crossing-exponent-dichotomy}.
No submultiplicativity of \(P_\kk\) is used.
\end{proof}

\begin{corollary}[The complex Hurwitz exponent has the same gap]
\label{cor:hurwitz-exponent-dichotomy}
Under the same inputs over \(\mathbb C\), the minimum
\(\sigma_\mathbb C(q)\) of complex bilinear squares in the
Hurwitz identity has exponent one if the square is not a price,
and exponent \(\sqrt2\) if it is a price.
\end{corollary}
\begin{proof}
The checked comparisons in Corollary~\ref{cor:square-hurwitz}
give
\(\sigma_\mathbb C(q)\leq2P_\mathbb C(q)\) and
\(P_\mathbb C(q)\leq2\sigma_\mathbb C(q)\).
The additive logarithmic errors vanish after division by \(\log q\).
\end{proof}

The near-linear bound on the nonprice branch is an existence
consequence of spectral duality. This proof supplies no explicit
sequence of maps attaining it and does not transfer the usual
sum-of-squares formulation to characteristic two.

\paragraph{The known noncommutative circuit reduction.}
Let \(C_\kk(f)\) denote circuit size counted by edges and let
\(\mathrm{PERM}_{4r}^{\mathrm{nc}}\) be the row-ordered noncommutative
permanent. The central width \(w\) of Hrube\v{s}, Wigderson and
Yehudayoff satisfies
\(w(\mathrm{ID}_q)=P_\kk(q)\), where
\(\mathrm{ID}_q=\sum_{i,j}x_i y_j x_i y_j\): a degree-four central
term has rank one across either \(12\mid34\) or \(14\mid23\), and
every such rank-one term has the corresponding central form.
Position filtering preserves these forms, and grouping the terms by
their cut identifies the minimum count with the minimum sum of ranks
\cite[Lemma~3.5, p.~883]{HrubesWigdersonYehudayoff2011}.
The source's circuit-to-width estimate and permanent substitution
\cite[Propositions~3.2 and~5.3, Lemma~5.5,
pp.~881--882, 888--889]{HrubesWigdersonYehudayoff2011}
therefore give, over every field,
\[
 C_\kk\bigl(\mathrm{PERM}_{4r}^{\mathrm{nc}}\bigr)
 \geq \Omega\!\left(\frac{P_\kk(2^r)}{r^3 2^r}\right).
\]
If the square is a price, this is at least
\(2^{(\sqrt2-1-o(1))r}\).
This conditional substitution into the known width reduction applies
also in characteristic two; it supplies no converse and concerns the
noncommutative permanent.
Liu's unrefereed Theorem~1.1 claims the stronger lower bound
\(2^{n/2-O(\sqrt{n\log n})}\) in order \(n\) under the absence of
a planar strict shortening \cite[p.~1]{Liu2026noncommutative}.
The role of the classical width reduction here is to show why even
a fixed superlinear lower exponent for \(P_\kk(q)\) entails a
substantial noncommutative circuit lower bound.
